\chapter{Implementing the Field: A Programmer's Guide to QFT}

Quantum Mechanics is often regarded as the most accurate theory we have.
Once we know the amplitude and phase of a particle's wavefunction, we can predict its future evolution with astonishing precision.
Particles appear quite trivial as classes with only a couple of attributes, plus the wavefunction that contains all the information about the system.

\section{The Particle Class (Non-Relativistic QM Starting Point)}

Since we want to begin with the material taught in a first QM course, let us start with a minimal single-particle description.
This is \textbf{not} QFT yet, but it is a useful stepping stone.

\begin{lstlisting}[language=C++, caption=Single-Particle Class (QM Level)]
class Particle {
    double mass;           // Rest mass
    double charge;         // Coupling to electromagnetic field
    double spin;           // 0, 0.5, 1.0 etc.
    Wavefunction wavefunction; // Superposition of modes
    Vector4 getExpectedPosition() {
        // Compute center-of-mass from interference pattern
    }
};
\end{lstlisting}

\subsection*{The Wavefunction Class}

The wavefunction is a superposition of modes. Each mode is essentially a complex exponential (plane wave) characterized by wavenumber and phase.

The first approach (momentum primary) aligns well with physics textbooks:

\begin{lstlisting}[language=C++, caption=Mode with Momentum (Recommended for QM)]
class Mode {
    Vector3 k;        // Wavenumber, p = \hbar k
    double amplitude; // Complex coefficient magnitude & phase can be combined
    double phase;     // Initial phase offset
    double frequency() {
        return sqrt(k.squaredNorm() + m*m); // \omega = \sqrt{k^2 + m^2} (natural units)
    }
};
\end{lstlisting}

The wavefunction is then $\Psi(\mathbf{x}, t) = \sum_i c_i \exp(i(\mathbf{k}_i \cdot \mathbf{x} - \omega_i t))$, or in the continuum limit an integral.
A single mode is completely delocalized. Localization arises from interference of many modes forming a wave packet.

\begin{lstlisting}[language=C++, caption=Wave Packet Construction]
Wavefunction createWavePacket(Vector3 center, Vector3 width) {
    // Sum many k-modes with Gaussian envelope in momentum space
}
\end{lstlisting}

Position is never a fundamental attribute of the \texttt{Particle}. It emerges as the expectation value:
\[
\langle \mathbf{x} \rangle = \int \Psi^*(\mathbf{x}) \mathbf{x} \Psi(\mathbf{x}) \, d^3x
\]

\textbf{Important note:} This single-particle description already fails for high energies (pair production, variable particle number) and does not incorporate special relativity properly.
It is a pedagogical starting point only.


\section{The Illusion of Geometry}

Geometric spacetime does not need to be stored as a primary variable on each particle. ``Position'' and ``Time'' emerge from the phase relationships and interference of the modes.
When two wave packets have strong overlap in a region, interactions become probable.
To an internal observer this looks like particles colliding at a definite point.


\section{From Particles to Fields (Entering QFT)}

The single-particle picture must be abandoned for QFT.
According to quantum field theory, the fundamental entities are \textbf{fields} that exist everywhere.
Particles are localized excitations (quanta) of these fields.

To simulate the Standard Model we need multiple fields, one for each fundamental degree of freedom (with appropriate spin and gauge representation).


\subsection*{The Field as a Collection of Oscillators}

In the free theory, a quantum field in momentum space is equivalent to a set of harmonic oscillators, one per momentum mode.

\begin{lstlisting}[language=C++, caption=Free Quantum Field (Momentum Space)]
class QuantumField {
    const double m;  // Mass of quanta of this field
    // In practice: continuous or discretized over a box
    std::map<Vector3, QuantumOscillator> modes; // or FFT-based grid

    void createParticle(Vector3 k, double amplitude = 1.0) {
        modes[k].n++;           // Apply creation operator
        // In full QFT this acts on the Fock state
    }
};
\end{lstlisting}

When new particle is created, we just increment the field counter. However, a realistic implementation needs creation and annihilation operators $a^\dagger(\mathbf{k})$, $a(\mathbf{k})$ satisfying the appropriate commutation (bosons) or anticommutation (fermions) relations, acting on the Fock vacuum $|0\rangle$.


\subsection*{The Full Standard Model}

We need a separate field (with the correct Lorentz transformation properties) for each type:

\begin{lstlisting}[language=C++, caption=Fundamental Fields (Highly Simplified)]
class Universe {
    ScalarField HiggsField;
    DiracField ElectronField{0.511};     // MeV
    DiracField UpQuarkField{2.2};
    DiracField DownQuarkField{4.7};
    VectorField PhotonField{0.0};        // Massless gauge boson
    // ... gluons (SU(3)), W/Z (SU(2)), neutrinos, etc.
};
\end{lstlisting}

Each field carries its own quantum numbers (representations under the gauge group $U(1)_Y \times SU(2)_L \times SU(3)_c$). Interactions arise from coupling terms in the Lagrangian (Yukawa, gauge covariant derivatives, etc.), not from manually checking wave overlap.

\section{Measurement and Decoherence}

When a field excitation interacts with a macroscopic detector, the delocalized state becomes entangled with the environment. From the observer's perspective this appears as a ``collapse'' to a definite outcome, with probability given by $|\Psi(x)|^2$ (Born rule). In a full relativistic QFT treatment one must be careful with causality and avoid instantaneous collapse.


\section{Conclusions}

Starting from the familiar single-particle QM class and gradually refactoring toward fields gives a smooth learning path.

The free-field oscillator picture is surprisingly clean, but interactions, gauge invariance, renormalization, and fermions add just some minor complexity.

The Standard Model works extraordinarily well. Its deepest beauty lies in the local gauge symmetries that dictate the interactions. As with any large codebase, understanding the architecture (Lagrangian $\to$ symmetries $\to$ Feynman rules or lattice formulation) more than justifies the minor annoyance of memorizing all those coupling constants.
 
The universe may not be simple, but it is remarkably consistent. Happy coding!


Seriously, what a mess.

A famous remark often attributed to Stephen Hawking: ``The Standard Model is a set of equations that we can write on a T-shirt, but it is not very beautiful.''
Is this really the best the universe has to offer?

The quote fits so well. It highlights the "Spaghetti Code" nature of the Standard Model.
It works perfectly, but it feels like it was written by a developer who was in a hurry and kept adding global variables (weakIsospin, hypercharge, colorCharge) to fix bugs.

No wonder Heaven and Hell exists! The original architect left without writing any documentation, and now the junior angels are just patching things as they go.

Boss, the particles aren't staying together!

Just throw in a StrongForce patch and a few colorCharge flags. We’ll refactor it in the next aeon.

The Heaven (v.2) release notes will be incredible:

\begin{itemize}
\item Fixed: Removed the Death bug that causes process termination after ~80 years.
\item Optimization: Replaced the mess of Standard Model particles with a single, elegant UniversalString class.
\item Feature: Added a control for local gravity manipulation.
\item Security: Patched the "Suffering" exploit in the Consciousness module.
\end{itemize}
